% !TEX root = ../main.tex
\chapteropening{%
\Cref{ch:geometry} gave us the left side of Einstein's equation, the geometry.
This chapter supplies the right side, the description of matter, and then puts
the two together.

By the end of it you will be able to take any metric, compute the matter it
requires, and read the result the way a spacetime engineer does: as a list of
demands that some physical material would have to meet.}{%
\item describe the energy, momentum and stress of matter with the
  stress-energy tensor, and say what each component means;
\item compute the energy density that any observer measures;
\item write down the stress-energy of dust, fluids, radiation, tension and
  vacuum energy;
\item state Einstein's equation, convert between geometric and everyday units,
  and recover Newton's gravity from it;
\item carry out the spacetime engineer's central calculation, from a chosen
  metric to the matter it demands, and recognize the first signs of an exotic
  demand.}
\chapter{How Matter Shapes Spacetime}\label{ch:matter-geometry}

% ===========================================================================
\section{Energy depends on who measures it}\label{sec:energy-observers}

\begin{thought}{A brick in passing}
A brick sits on a table. You fly past it at 80 percent of the speed of light.
Because the brick moves relative to you, each of its atoms carries more energy
by the factor $\gamma = 1/\sqrt{1 - 0.8^{2}} = 5/3$. Is the energy per unit
volume that you measure larger by the same factor, or by a different one?
\end{thought}

In Newtonian physics, the mass of a brick is one fixed number. In relativity
the brick's \emph{energy} depends on the observer: someone moving past it sees
it carry kinetic energy as well as its rest energy. To describe matter in a
way every observer can use, we need to say how the measured energy changes
from one observer to another.

A particle of mass $m$ with four-velocity $u^{\mu}$ has \emph{four-momentum}
$p^{\mu} = m\,u^{\mu}$. An observer with four-velocity $w^{\mu}$ measures the
particle's energy to be
\begin{equation}
  E = -g_{\mu\nu}\,p^{\mu}w^{\nu} \equiv -p_{\mu}w^{\mu}.
  \label{eq:measured-energy}
\end{equation}
In the second form we have \emph{lowered an index}, $p_{\mu} = g_{\mu\nu}p^{\nu}$:
lowering indices with the metric is the bookkeeping that turns a pair of
vectors into a number every observer agrees on. For an observer at rest in flat
spacetime, $w^{\mu} = (1,0,0,0)$ and \cref{eq:measured-energy} gives
$E = p^{0}$. For a particle at rest relative to the observer, $E = m$, its rest
energy $mc^{2}$ in ordinary units.

A continuous material, such as a gas, a fluid or a stretched rope, needs more
than a four-momentum. At each event we want to know how much energy is present
per unit volume, how fast energy is flowing, how much momentum there is, and
what forces neighbouring parts of the material exert on each other. All of
this fits into one symmetric $4\times4$ array, the \emph{stress-energy tensor}
$T^{\mu\nu}$ (\cref{fig:stress-energy}). In the frame of a given observer,
$T^{tt}$ is the \textbf{energy density}; $T^{ti} = T^{it}$ is the
\textbf{energy flux} in direction $i$, which equals the density of momentum in
that direction; and $T^{ij}$ is the \textbf{stress}, the flux of $i$-momentum
across a surface facing the $j$ direction. The diagonal stresses are
\textbf{pressures} and the off-diagonal ones are \textbf{shear stresses}.

\begin{figure}[tb]
\centering
\includegraphics{ch03-stress-energy}
\caption{The components of the stress-energy tensor in one observer's frame.
The array is symmetric, so ten numbers describe the matter at each event.}
\label{fig:stress-energy}
\end{figure}

Observers in relative motion split the same tensor into energy, flux and
stress differently, just as they split an interval into time and distance
differently. The rule generalizes \cref{eq:measured-energy}: an observer with
four-velocity $w^{\mu}$ measures the energy density
\begin{equation}
  \rho_{w} = T_{\mu\nu}\,w^{\mu}w^{\nu},
  \label{eq:observed-density}
\end{equation}
where $T_{\mu\nu} = g_{\mu\alpha}g_{\nu\beta}T^{\alpha\beta}$ is the tensor
with both indices lowered. This formula is one of the most used in the book.
When a geometry demands matter, the first question is always what energy
density each observer would measure.

The simplest material is \textbf{dust}: particles at rest relative to one
another, with no pressure between them. In the dust's own rest frame the only
nonzero component is the energy density $\rho_{0}$, and the tensor that
reduces to that in the rest frame is
\begin{equation}
  T^{\mu\nu} = \rho_{0}\,u^{\mu}u^{\nu},
  \label{eq:dust}
\end{equation}
where $u^{\mu}$ is the dust's four-velocity.

\begin{example}[label=ex:tension-passing]{Tension seen in passing}
A long rod at rest has energy density $\rho_{0}$ and is stretched along its
length, the $x$ direction, by a tension $\tau$. In its rest frame
$T^{tt} = \rho_{0}$ and $T^{xx} = -\tau$ (a pull is a negative pressure), and
all other components vanish. What energy density does an observer moving
along the rod at speed $v$ measure?

\solution The observer's four-velocity is $w^{\mu} = \gamma\,(1, v, 0, 0)$. In
flat spacetime, lowering both indices leaves these two components unchanged:
$T_{tt} = \rho_{0}$ and $T_{xx} = -\tau$. \Cref{eq:observed-density} then gives
\begin{equation*}
  \rho_{w} = \gamma^{2}\bigl(T_{tt} + v^{2}\,T_{xx}\bigr)
  = \gamma^{2}\bigl(\rho_{0} - v^{2}\tau\bigr).
\end{equation*}

\discussion Tension lowers the energy density that a passing observer
measures. For ordinary materials the effect is tiny: the energy density
includes the rest energy, and for steel near its breaking point $\tau/\rho_{0}$
is about $10^{-12}$. But a material whose tension \emph{exceeded} its energy
density would look negative to observers moving fast enough along it, since
$\rho_{w} < 0$ whenever $v^{2} > \rho_{0}/\tau$. That is exactly the kind of
material Morris and Thorne found a wormhole throat needs.
\end{example}

\begin{checkpoint}[label=chk:laser]
A laser beam travelling along $+x$ has $T^{tt} = T^{tx} = T^{xt} = T^{xx} = u$,
with all other components zero. What energy density does an observer moving
along $x$ at speed $v$ measure? What does one moving at $0.8$ with the beam
measure, and one moving at $0.8$ against it?
\end{checkpoint}

\begin{revisited}{A brick in passing}
By more: the factor is $\gamma^{2}$. Each atom's energy is larger by $\gamma$,
and the brick is contracted along the direction of motion, which packs the
atoms more densely by another factor of $\gamma$. The formulas say the same
thing: with $T^{\mu\nu} = \rho_{0}u^{\mu}u^{\nu}$, \cref{eq:observed-density}
gives $\rho_{w} = \rho_{0}(u_{\mu}w^{\mu})^{2} = \gamma^{2}\rho_{0}$, because
$u_{\mu}w^{\mu} = -\gamma$ for two four-velocities in relative motion at speed
$v$. At 80 percent of the speed of light, $\gamma^{2} = 25/9 \approx 2.8$.
Energy density is not a quantity all observers share; it is one component of a
tensor, and each observer sees a different mix of components.
\end{revisited}

% ===========================================================================
\section{Pressure, tension and the vacuum}\label{sec:catalogue}

\begin{thought}{Flying through empty space}
Cosmologists have found that empty space seems to carry a small energy density
of its own, the \enquote{dark energy} that makes the expansion of the universe
speed up. When you flew past the brick, you measured more energy density than
someone at rest beside it. If you fly through empty space at 80 percent of the
speed of light, do you measure more of the vacuum's energy density too?
\end{thought}

The stress-energy tensors of a few kinds of matter recur throughout this book.
Each is simplest in its own rest frame.

\paragraph{Perfect fluid.} A fluid with energy density $\rho$ and the same
pressure $p$ in every direction has
\begin{equation}
  T^{\mu\nu} = (\rho + p)\,u^{\mu}u^{\nu} + p\,g^{\mu\nu},
  \label{eq:perfect-fluid}
\end{equation}
which in the rest frame is $\mathrm{diag}(\rho, p, p, p)$. Dust is a perfect fluid
with $p = 0$. Gases, liquids and the interiors of stars are perfect fluids to
a good approximation.

\paragraph{Radiation.} A gas of light moving in all directions is a perfect
fluid with $p = \rho/3$.

\paragraph{Tension.} Stretch a rope and each part pulls on its neighbours along
the rope. A pull is a negative pressure, so tension appears in $T^{\mu\nu}$ as a
\emph{negative} diagonal stress: a material under tension $\tau$ along $z$ has
$T^{zz} = -\tau$. Tension is one of the spacetime engineer's most important
ingredients, and you will meet it again and again in the demands of engineered
geometries.

\paragraph{Vacuum energy.} Empty space can carry an energy density
$\rho_{\text{vac}}$ of its own, with stress-energy
\begin{equation}
  T_{\mu\nu} = -\rho_{\text{vac}}\,g_{\mu\nu}.
  \label{eq:vacuum-energy}
\end{equation}
In its rest frame, if it can be said to have one, this is an energy density
$\rho_{\text{vac}}$ with pressure $p = -\rho_{\text{vac}}$ in every direction: a
tension equal to the energy density.

\begin{deeper}{the electromagnetic field}
In Heaviside--Lorentz units with $c = 1$, a magnetic field of strength $B$
along $z$ has energy density $B^{2}/2$, a pressure $+B^{2}/2$ across the field
lines ($T^{xx} = T^{yy}$), and a tension $B^{2}/2$ along them
($T^{zz} = -B^{2}/2$). Field lines behave like stretched elastic bands that
push sideways on one another. The tension along the field equals the energy
density: this is the largest tension the dominant energy condition of
Chapter~5 allows, and magnetic fields reach it exactly.
\end{deeper}

\begin{checkpoint}[label=chk:rho-p]
For dust, radiation, a perfect fluid with $p > 0$, and vacuum energy, compute
$\rho + p$. Which of them has $\rho + p = 0$? (In Chapter~5 you will see that
the sign of $\rho + p$ decides whether matter focuses light.)
\end{checkpoint}

\begin{revisited}{Flying through empty space}
No: every observer measures the same vacuum energy density. For any
four-velocity, \cref{eq:vacuum-energy} and \cref{eq:four-velocity} give
$\rho_{w} = -\rho_{\text{vac}}\,g_{\mu\nu}w^{\mu}w^{\nu} = \rho_{\text{vac}}$.
In the language of \cref{ex:tension-passing}, the vacuum's tension equals its
energy density, so the factor $\gamma^{2}(\rho - v^{2}\tau)$ becomes
$\gamma^{2}\rho_{\text{vac}}(1 - v^{2}) = \rho_{\text{vac}}$: the tension exactly
cancels the boost that made the brick look denser. That is why a vacuum energy
defines no preferred state of rest, as empty space should not.
\end{revisited}

% ===========================================================================
\section{Einstein's equation}\label{sec:einstein}

\begin{thought}{Between the Earth and the Moon}
Between the Earth and the Moon there is almost no matter at all, so the
stress-energy on the right side of Einstein's equation is zero there. Does the
equation then say that spacetime between them is flat? If it does, what steers
the Moon around its orbit, and what raises the tides in the Earth's oceans?
\end{thought}

Energy and momentum are conserved: whatever flows out of a small box must be
accounted for by a loss inside it. In flat spacetime this is the statement
\begin{equation}
  \partial_{\mu}T^{\mu\nu} = 0 .
  \label{eq:conservation-flat}
\end{equation}
Its $\nu = t$ component is the continuity equation for energy,
$\partial_{t}T^{tt} + \partial_{i}T^{it} = 0$; its three spatial components say
that momentum changes only through the stresses acting on a region. In curved
spacetime the partial derivative becomes the \emph{covariant} derivative
$\nabla_{\mu}$, which adds Christoffel-symbol terms: $\nabla_{\mu}T^{\mu\nu} = 0$.
Physically, matter then exchanges energy and momentum with the gravitational
field in a definite way, and the extra terms describe the exchange.

Einstein's equation is
\begin{equation}
  G_{\mu\nu} = \frac{8\pi G}{c^{4}}\,T_{\mu\nu},
  \label{eq:einstein-si}
\end{equation}
with $G_{\mu\nu}$ the Einstein tensor of \cref{eq:einstein-tensor}. There are ten
equations here, one for each independent component of the two symmetric
tensors.

\begin{deeper}{why the Einstein tensor?}
Whatever sits on the left side must be consistent with conservation of energy
and momentum on the right. The Einstein tensor has exactly this property: its
covariant divergence vanishes identically, $\nabla^{\mu}G_{\mu\nu} = 0$ for
\emph{every} metric, a result called the contracted Bianchi identity.
Einstein's equation therefore implies $\nabla^{\mu}T_{\mu\nu} = 0$
automatically. Apart from a term proportional to the metric itself, which acts
as vacuum energy, the Einstein tensor is essentially the only combination of
the metric and its first two derivatives with this property.
\end{deeper}

From now on we set $G = 1$ as well as $c = 1$, and Einstein's equation reads
\begin{equation}
  G_{\mu\nu} = 8\pi\,T_{\mu\nu}.
  \label{eq:einstein}
\end{equation}
In these \emph{geometric units} everything is measured in powers of length. A
mass $M$ becomes the length $GM/c^{2}$: the Sun is \SI{1.48}{\kilo\metre} and
the Earth \SI{4.4}{\milli\metre}. Curvature has units of one over length
squared, and so, by \cref{eq:einstein}, do energy density and stress. To
convert back, multiply a geometric energy density or stress, with lengths in
metres, by
\begin{equation}
  \frac{c^{4}}{G} = \SI{1.21e44}{\newton}
  \label{eq:conversion}
\end{equation}
to get joules per cubic metre, which are the same as pascals, and divide by
$c^{2}$ to express an energy density as a mass density. This one conversion
factor is the stiffness of spacetime from \cref{eq:stiffness}.

Einstein's equation also has a statement in plain words, due to John Baez and
Emory Bunn. Release a small ball of test particles, all at rest relative to
one another, and let them fall freely. Tides may at once begin to change the
ball's shape, and its volume $V$ begins to change at the rate
\begin{equation}
  \frac{\ddot V}{V} = -4\pi\,\bigl(\rho + p_{x} + p_{y} + p_{z}\bigr),
  \label{eq:ball}
\end{equation}
where the dots are derivatives with respect to the particles' proper time, and
the energy density and the three pressures are those at the centre of the
ball, measured by an observer at rest relative to it. Energy density makes the
ball shrink, and so does pressure. Required of every small ball, in every
state of motion, \cref{eq:ball} is equivalent to the whole of Einstein's
equation.

\begin{example}[label=ex:newtonian, unbreakable]{The Newtonian limit}
Show that for the weak-field metric \cref{eq:weak-field} with a static
potential $\Phi$, the $tt$ component of Einstein's equation reduces to Poisson's
equation of Newtonian gravity, $\nabla^{2}\Phi = 4\pi\rho$.

\solution Compute the Einstein tensor of \cref{eq:weak-field} to first order in
$\Phi$. The Christoffel symbols are first order, so their products in
\cref{eq:riemann} are second order and can be dropped. The calculation
(\cref{prob:weak-field}) gives $G_{tt} \approx 2\,\nabla^{2}\Phi$. For slowly
moving matter the dominant component of the stress-energy is the energy
density, $T_{tt} \approx \rho$, so $G_{tt} = 8\pi T_{tt}$ becomes
$\nabla^{2}\Phi = 4\pi\rho$: Poisson's equation, in units with $G = 1$.

\discussion Newtonian gravity is the weak-field, slow-motion limit of
Einstein's equation, and the factor $8\pi$ is what makes the limit come out
right. Einstein fixed the constant in just this way, by requiring his equation
to reproduce Newton's gravity wherever Newton's gravity works. Every
prediction beyond Newton, from the bending of starlight to the clock rates of
\cref{ch:geometry}, then came with no adjustable number left.
\end{example}

In empty space the right side of \cref{eq:ball} vanishes. A ball of falling
particles is still stretched and squeezed by tides, but its volume does not
begin to change: the stretching along one direction is balanced by squeezing
along the others, as you found for the Earth in \cref{chk:tidal-volume}. In
Newtonian terms, Poisson's equation in empty space says $\nabla^{2}\Phi = 0$,
not $\Phi = 0$. In the full theory, the trace of Einstein's equation gives
$R = -8\pi T$, with $T = g^{\mu\nu}T_{\mu\nu}$, so in empty space $R = 0$ and
then $G_{\mu\nu} = 0$ says $R_{\mu\nu} = 0$. The Ricci tensor, which governs
changes of volume, vanishes. The Riemann tensor, with twenty independent
components to the Ricci tensor's ten, need not; the rest of the curvature is
set by matter elsewhere and by what arrives from far away.

\begin{checkpoint}[label=chk:water]
Water has a density of \SI{1000}{\kilogram\per\cubic\metre}. What is its
energy density in geometric units? Einstein's equation then sets a curvature
of order $8\pi\rho$; over what distance does spacetime filled with water curve
appreciably?
\end{checkpoint}

\begin{revisited}{Between the Earth and the Moon}
No. Einstein's equation fixes only part of the curvature. Where there is no
matter it sets the Ricci tensor to zero, so a ball of falling particles keeps
its volume at first, but the rest of the Riemann tensor survives. In Newtonian
terms, the potential between the Earth and the Moon obeys $\nabla^{2}\Phi = 0$
without being zero: its slope steers the Moon around its orbit, and its second
derivatives, which add up to zero, stretch the oceans along the line to the
Moon and squeeze them across it. Empty space can be curved. Matter is what
makes a falling ball begin to shrink.
\end{revisited}

% ===========================================================================
\section{Pressure gravitates}\label{sec:pressure}

\begin{thought}{A box of hot gas}
Inside a star, pressure adds to the source of gravity. Take a sealed steel box
of gas and heat it until the pressure inside is enormous, while the box keeps
its shape. Does the box now pull on a distant planet harder than its energy
alone would suggest, because of all that pressure?
\end{thought}

Einstein's equation says more than Newton's. In a perfect fluid at rest the
three pressures are equal, so by \cref{eq:ball} a small ball of test particles
begins to shrink at a rate set by
\begin{equation}
  \rho + 3p ,
  \label{eq:active-mass}
\end{equation}
and this is the quantity that plays the role of $\rho$ in Poisson's equation
(\cref{prob:rho-3p}).
Pressure is a source of gravity alongside energy. For ordinary matter the
pressure is tiny compared with the energy density, which includes the rest
energy $mc^{2}$, so Newton never noticed. Tension, however, \emph{weakens}
gravity, and enough of it reverses its sign. Vacuum energy has
$\rho + 3p = -2\rho_{\text{vac}}$: it repels. This is why a positive vacuum
energy makes the expansion of the universe accelerate, and it is a first hint
of what engineered spacetimes can do with tension.

\begin{keyidea}
In general relativity every component of the stress-energy tensor bends
spacetime: energy, momentum and stress alike. Pressure attracts, tension
repels, and a spacetime engineer designs with all of them.
\end{keyidea}

\begin{checkpoint}[label=chk:stop-tension]
A material has energy density $\rho$ and the same stress $p$ in every
direction. What value of $p$ would keep a small ball of test particles inside
it from beginning to shrink? Is that a pressure or a tension?
\end{checkpoint}

\begin{figure}[tb]
\centering
\includegraphics{ch03-gas-box}
\caption{A sealed box of hot gas, cut by an imaginary plane. Across the cut,
the gas pushes the two halves apart and the stretched walls pull them
together. In equilibrium the forces balance, for every cut.}
\label{fig:gas-box}
\end{figure}

\begin{revisited}{A box of hot gas}
No. The gas pressure does add $3p$ to the source of gravity inside the box,
but the steel walls that hold the gas in are under tension, and tension
subtracts. For any object that holds itself together in equilibrium, the two
cancel exactly. Cut the box in two with an imaginary plane
(\cref{fig:gas-box}). The gas pushes the halves apart with a force equal to its
pressure times the area of gas in the cut, and the walls pull them together
with their tension times the area of steel in the cut. Equilibrium means these
balance, for every cut in every direction, so the total of all the stresses
over the box, $\int (T^{xx} + T^{yy} + T^{zz})\,\dd V$, is zero
(\cref{prob:von-laue}). A distant planet feels only the box's total energy,
which heating has raised by the thermal energy added, $\Delta E$, making the
box heavier by $\Delta E/c^{2}$.
\end{revisited}

% ===========================================================================
\section{Reading the equation backwards}\label{sec:backwards}

\begin{thought}{A room where time runs fast}
In a rocket accelerating through empty space, clocks at the nose tick faster
than clocks at the tail, and nothing but the rocket's acceleration is needed to
make it so. Could you instead arrange for clocks to tick fastest in the middle
of a room and slower toward both walls, keeping the room's space perfectly
flat? What would have to fill the room?
\end{thought}

Here is the spacetime engineer's central calculation. It has five steps.
\begin{enumerate}
\item Write down the metric $g_{\mu\nu}$ of the geometry you want.
\item Compute the Christoffel symbols, \cref{eq:christoffel}.
\item Compute the Riemann tensor \cref{eq:riemann}, then the Ricci tensor, the
  Ricci scalar and the Einstein tensor \cref{eq:einstein-tensor}.
\item Read off the demanded stress-energy, $T_{\mu\nu} = G_{\mu\nu}/8\pi$.
\item Interpret it: compute the energy density and stresses measured by the
  observers you care about, using \cref{eq:observed-density}.
\end{enumerate}
For all but the simplest metrics, steps 2 and 3 are best done with computer
algebra. The examples in this chapter are simple enough to follow by hand, and
each teaches something the rest of the book builds on.

\begin{example}[label=ex:lapse, unbreakable]{Clocks that run at different rates}
Suppose we want clocks to run at different rates at different places along
$x$, while space itself stays exactly flat. The metric is
\begin{equation}
  \dd s^{2} = -\alpha(x)^{2}\,\dd t^{2} + \dd x^{2} + \dd y^{2} + \dd z^{2},
  \label{eq:lapse-metric}
\end{equation}
where $\alpha(x) > 0$ is the rate of a clock at rest at $x$, relative to
coordinate time. Find the matter this geometry demands.

\solution \emph{Christoffel symbols.} Only $g_{tt} = -\alpha^{2}$ varies, and
only with $x$. Writing $\alpha' = \dd\alpha/\dd x$, \cref{eq:christoffel} gives
two independent nonzero symbols:
\begin{equation*}
  \Gamma^{t}{}_{tx} = \Gamma^{t}{}_{xt} = \frac{\alpha'}{\alpha},
  \qquad
  \Gamma^{x}{}_{tt} = \alpha\,\alpha' .
\end{equation*}
\emph{Curvature.} From \cref{eq:riemann}, the nonzero components of the Ricci
tensor are
\begin{equation*}
  R_{tt} = \alpha\,\alpha'', \qquad R_{xx} = -\frac{\alpha''}{\alpha},
\end{equation*}
and the Ricci scalar is $R = g^{tt}R_{tt} + g^{xx}R_{xx} = -2\alpha''/\alpha$.
\emph{Einstein tensor.} From \cref{eq:einstein-tensor}, $G_{tt} = 0$ and
$G_{xx} = 0$, while
\begin{equation*}
  G_{yy} = G_{zz} = -\tfrac{1}{2}R = \frac{\alpha''}{\alpha}.
\end{equation*}
\emph{Interpretation.} A clock at rest at $x$ has four-velocity
$w^{\mu} = (1/\alpha, 0, 0, 0)$, and by \cref{eq:observed-density} it measures
energy density $\rho = G_{tt}/(8\pi\alpha^{2}) = 0$. The demand is
\begin{equation}
  \rho = 0, \quad p_{x} = 0, \quad p_{y} = p_{z} = \frac{\alpha''}{8\pi\,\alpha}.
  \label{eq:lapse-demand}
\end{equation}

\discussion Changing the rate of clocks from place to place, with space kept
flat, demands \emph{no energy at all}: only stress, and only in the directions
perpendicular to the change. If $\alpha$ is a straight line, $\alpha = 1 + gx$,
the demand vanishes entirely: that geometry is flat spacetime seen by an
observer with constant acceleration $g$ (\cref{prob:rindler}), the rocket of
the thought experiment.
\end{example}

The Earth also makes clocks run at different rates at different heights, yet
the empty space around it holds no stress. The difference is that the Earth's
geometry curves space as well: the factor $(1 - 2\Phi)$ in front of the
spatial part of \cref{eq:weak-field}. In empty space the curvature of space and
the variation of clock rates cancel exactly in Einstein's equation
(\cref{prob:weak-field}). In \cref{ex:lapse} we insisted on keeping space flat,
and the stress is the price of that choice. Designs make choices like this all
the time, and Einstein's equation prices each one.

The price is steep. Take a clock rate that changes by one part in a billion
across one metre, so that $\alpha'' \sim \SI{e-9}{\per\metre\squared}$. By
\cref{eq:lapse-demand,eq:conversion}, the stress is about
$10^{-9}/(8\pi) \times \SI{1.2e44}{\pascal} \approx \SI{5e33}{\pascal}$, the
pressure deep inside a neutron star, demanded to make two clocks a metre apart
differ by one part in a billion.

\begin{figure}[tb]
\centering
\includegraphics{ch03-lapse-bump}
\caption{Clocks that run fastest in the middle, $\alpha = 1 + \tfrac12
e^{-x^{2}}$ (top), and the stress this demands across the direction of the
change, \cref{eq:lapse-demand} (bottom). Around the peak the demand is a
tension; on the flanks it is a pressure. The energy density is zero
everywhere.}
\label{fig:lapse-bump}
\end{figure}

\begin{checkpoint}[label=chk:lapse-dip]
Suppose instead that clocks run \emph{slowest} in the middle, with
$\alpha = 1 - \tfrac12 e^{-x^{2}}$. Where is the demand a tension, and where a
pressure?
\end{checkpoint}

\begin{revisited}{A room where time runs fast}
The room must hold a strange material: one with no energy density at all, but
with stress along the two directions parallel to the walls
(\cref{fig:lapse-bump}). Where the clock rate peaks, $\alpha'' < 0$ and the
stress is a tension; nearer the walls, where $\alpha$ curves back up, it is a
pressure. The rocket needed nothing because its $\alpha$ is a straight line;
the room pays for the curvature of $\alpha$. And the tension at the peak is
worse than it looks. For light travelling parallel to the walls,
$\rho + p_{y} = \alpha''/(8\pi\alpha) < 0$, and in Chapter~5 you will learn that
ordinary matter always has $\rho + p \ge 0$ along every direction light can
travel. A simple bump in the rate of clocks already demands matter that no
ordinary material can provide.
\end{revisited}

% ===========================================================================
\section{When space flows unevenly}\label{sec:shear}

\begin{thought}{A river with a fast lane}
In \cref{ch:geometry}, space flowing at one uniform speed turned out to be flat
spacetime in disguise, and it cost nothing. Now let the flow run faster down
the middle of a channel than near its banks, the way a river does. Does that
cost anything? If it does, will the observers carried along by the flow
measure a positive or a negative energy density? Make a guess before reading
on.
\end{thought}

\begin{example}[label=ex:shear]{A flow of space that shears}
Let the speed of the flow of \cref{ex:flow} vary across space:
\begin{equation}
  \dd s^{2} = -\dd t^{2} + \bigl(\dd x + \beta(y)\,\dd t\bigr)^{2} + \dd y^{2} + \dd z^{2},
  \label{eq:shear-metric}
\end{equation}
a flow of space along $x$ whose speed depends on $y$. Find the energy density
measured by observers carried along by the flow.

\solution Observers carried by the flow move with $\dd x/\dd t = -\beta$ and
feel no force. Their four-velocity is $w^{\mu} = (1, -\beta, 0, 0)$, which
satisfies $g_{\mu\nu}w^{\mu}w^{\nu} = -1$ (\cref{prob:shear-check}). Carrying
out the recipe for \cref{eq:shear-metric} is longer than for \cref{ex:lapse},
and a good exercise for computer algebra (\cref{prob:shear-check}). The energy
density these observers measure is
\begin{equation}
  \rho = T_{\mu\nu}w^{\mu}w^{\nu} = -\frac{1}{32\pi}\left(\frac{\dd\beta}{\dd y}\right)^{2}.
  \label{eq:shear-density}
\end{equation}

\discussion Where the flow is uniform, $\dd\beta/\dd y = 0$ and nothing is
demanded, just as in \cref{ex:flow}. Wherever the flow shears, the observers
riding it measure a \emph{negative} energy density, whichever way the shear
goes, since the derivative is squared.
\end{example}

The numbers are once again enormous. A flow that changes from zero to the speed
of light across one metre has $\dd\beta/\dd y \approx \SI{1}{\per\metre}$, and
\cref{eq:shear-density,eq:conversion} give an energy density of about
\SI{-1.2e42}{\joule\per\cubic\metre}, equivalent to a mass density of about
\SI{-1e25}{\kilogram\per\cubic\metre}. Nuclear matter, the densest material in
nature, has about \SI{+2e17}{\kilogram\per\cubic\metre}.

\begin{checkpoint}[label=chk:shear-width]
Suppose the flow speed in \cref{ex:shear} rises steadily from $0$ to
$\beta_{0}$ across a layer of width $w$. What energy density do the riding
observers measure in the layer? If you make the layer twice as wide, what
happens to the energy density, and to the total negative energy in the layer
per unit area of its face?
\end{checkpoint}

\begin{revisited}{A river with a fast lane}
It costs a great deal, and the cost is negative energy. By
\cref{eq:shear-density}, observers riding a shearing flow of space measure
$\rho = -(\dd\beta/\dd y)^{2}/32\pi$, negative whichever way the shear runs.
This is the warp drive's problem in miniature. At the wall of a warp bubble the
flow of space changes from the bubble's speed to zero, so the wall is a region
of shear, and \cref{eq:shear-density} is why Alcubierre's bubble demands
negative energy there (\cref{fig:warp}b).
\end{revisited}

% ===========================================================================
\section{Every geometry has a price}\label{sec:price}

The recipe of \topicref{sec:backwards} never fails. Every metric you can write
down has an Einstein tensor, so every geometry has a demanded stress-energy.
Read backwards, Einstein's equation is a price list: it tells you exactly what
matter any geometry would require.

Whether anything in nature can pay the price is a separate question, and it is
the question the rest of this book is organized around. The examples of this
chapter already show its shape. Shaping the rate of clocks demands stress
without energy, and can demand matter that violates $\rho + p \ge 0$. Shearing
the flow of space demands negative energy density. Chapter~4 splits spacetime
into space and time the way these examples suggest, with a \emph{lapse} that
sets the rate of clocks and a \emph{shift} that describes the flow of space, and
shows how to compute demands like these efficiently. Chapter~5 states precisely
which demands ordinary matter can meet, and Chapter~6 asks how far quantum
physics can go beyond them.

% ===========================================================================
\begin{chaptersummary}
\item Energy depends on the observer. An observer with four-velocity $w^{\mu}$
  measures the energy density $\rho_{w} = T_{\mu\nu}w^{\mu}w^{\nu}$; a passing
  observer sees a brick's energy density raised by $\gamma^{2}$.
\item The stress-energy tensor $T^{\mu\nu}$ collects energy density, energy
  flux, momentum density and stress into one symmetric array.
\item Dust, perfect fluids, radiation, tension and vacuum energy each have a
  characteristic stress-energy. Tension is negative pressure, and it lowers the
  energy density that passing observers measure.
\item Einstein's equation $G_{\mu\nu} = 8\pi T_{\mu\nu}$ (with $G = c = 1$) links
  curvature to matter. Multiply geometric stresses by
  $c^{4}/G = \SI{1.21e44}{\newton}$ to get pascals.
\item In plain words, Einstein's equation says that a small ball of test
  particles, initially at rest relative to one another, begins to change
  volume at the rate $\ddot V/V = -4\pi(\rho + p_{x} + p_{y} + p_{z})$. In empty
  space $R_{\mu\nu} = 0$: tides remain, but they change only the ball's shape.
\item Newton's gravity is the weak-field limit. Beyond it, pressure gravitates
  through $\rho + 3p$ and tension can make gravity repel; in a body that holds
  itself together, the pressures and tensions cancel in total.
\item Reading the equation backwards gives the matter any geometry demands.
  Varying clock rates across flat space demands stress without energy; a
  shearing flow of space demands negative energy density.
\end{chaptersummary}

\begin{checkanswers}
\item[\ref{chk:laser}] Lowering the indices flips the sign of each component
  with a single $t$: $T_{tt} = T_{xx} = u$ and $T_{tx} = T_{xt} = -u$. With
  $w^{\mu} = \gamma\,(1, v, 0, 0)$,
  $\rho_{w} = \gamma^{2}u\,(1 - 2v + v^{2}) = u\,(1 - v)/(1 + v)$. Chasing the
  beam at $0.8$ gives $u/9$; running into it ($v = -0.8$) gives $9u$. Each
  photon is Doppler shifted by the factor $\sqrt{(1 - v)/(1 + v)}$, and the
  photons' number density changes by the same factor.
\item[\ref{chk:rho-p}] Dust: $\rho + p = \rho$. Radiation: $\rho + p = \tfrac43\rho$.
  Fluid with $p > 0$: $\rho + p > \rho$. Vacuum energy: $\rho + p = 0$. Only
  vacuum energy sits exactly on the boundary.
\item[\ref{chk:water}] $\rho = G\rho_{\mathrm{mass}}/c^{2} = 6.67\times10^{-11}
  \times 1000/(9.0\times10^{16}) \approx \SI{7.4e-25}{\per\square\metre}$. A
  curvature of $8\pi\rho$ corresponds to a length
  $1/\sqrt{8\pi\rho} \approx \SI{2e11}{\metre}$, more than the distance from
  the Earth to the Sun. Spacetime is stiff.
\item[\ref{chk:stop-tension}] By \cref{eq:ball} the ball keeps its volume when
  $\rho + 3p = 0$, that is, $p = -\rho/3$: a tension of one third of the
  energy density in every direction. Vacuum energy, with $p = -\rho$, goes
  beyond this, and a ball of test particles in it begins to grow.
\item[\ref{chk:lapse-dip}] Now $\alpha'' > 0$ near $x = 0$, so the demand there
  is a pressure; on the flanks, where $\alpha$ curves back toward 1 from
  below, it is a tension. Wherever the demand is a tension, $\rho + p_{y} < 0$.
\item[\ref{chk:shear-width}] In the layer $\dd\beta/\dd y = \beta_{0}/w$, so
  $\rho = -\beta_{0}^{2}/(32\pi w^{2})$. Doubling $w$ divides the energy
  density by four. The layer then holds twice the volume, so the total per
  unit area, $-\beta_{0}^{2}/(32\pi w)$, is halved. Thin layers of shear are
  expensive on both counts.
\end{checkanswers}

\begin{problems}
\item \diff{1} A cloud of dust with rest-frame density $\rho_{0}$ moves at speed
  $v$ along $y$. Write down all components of $T^{\mu\nu}$ in the observer's
  frame.
\item \diff{1} \textbf{Geometric units.} Express the mass of the Earth, the Sun
  and a $10^{9}\,\Msun$ black hole in metres. What energy density, in joules per
  cubic metre, corresponds to a geometric energy density of
  \SI{1}{\per\square\kilo\metre}?
\item \diff{2} \textbf{Magnetic tension.} Using the stress-energy of a magnetic
  field in the \emph{Going deeper} box on the electromagnetic field, verify
  that $\rho + p \ge 0$ for every direction of the pressure, and that the tension
  along the field lines equals the energy density.
\item \diff{2} \textbf{Vacuum energy.} Show from \cref{eq:vacuum-energy} that
  every observer measures the same energy density $\rho_{\text{vac}}$, and that
  the pressure is $-\rho_{\text{vac}}$ in every direction.
\item \diff{2}\label{prob:rho-3p} \textbf{Pressure gravitates.} Einstein's
  equation can be rewritten as
  $R_{\mu\nu} = 8\pi\bigl(T_{\mu\nu} - \tfrac12 T g_{\mu\nu}\bigr)$, with
  $T = g^{\mu\nu}T_{\mu\nu}$. For a perfect fluid at rest, show that
  $R_{tt} = 4\pi(\rho + 3p)$. Given that $R_{tt} \approx \nabla^{2}\Phi$ in a weak,
  static field, explain why $\rho + 3p$ plays the role of $\rho$ in Poisson's
  equation.
\item \diff{3}\label{prob:von-laue} \textbf{Stresses in a body at rest.} For a
  static body of finite size in flat spacetime, conservation gives
  $\partial_{j}T^{ij} = 0$. Show that $\int T^{xx}\,\dd V = 0$ by integrating
  $\partial_{j}(x\,T^{xj}) = T^{xx}$ over a region enclosing the body, and deduce
  that the pressures and tensions in the box of hot gas cancel in total.
\item \diff{2}\label{prob:weak-field} \textbf{Weak fields.} Compute the Einstein
  tensor of \cref{eq:weak-field} to first order in $\Phi$. Show that
  $G_{tt} \approx 2\nabla^{2}\Phi$, and that all components vanish where
  $\nabla^{2}\Phi = 0$. Explain why the result differs from \cref{ex:lapse}.
\item \diff{2}\label{prob:rindler} \textbf{Uniform acceleration.} Show that the
  metric \cref{eq:lapse-metric} with $\alpha = 1 + gx$ has vanishing Riemann
  tensor. What does an observer at rest at fixed $x$ feel?
\item \diff{2} \textbf{A rippled clock rate.} Take $\alpha = 1 + \epsilon\sin kx$
  with $0 < \epsilon < 1$ in \cref{ex:lapse}. Where is the demand a tension? For
  light moving along $y$, where is $\rho + p_{y}$ negative?
\item \diff{3}\label{prob:shear-check} \textbf{The shearing flow.} Using a
  computer algebra system, compute the Einstein tensor of \cref{eq:shear-metric}
  and verify \cref{eq:shear-density}. Check that $w^{\mu} = (1,-\beta,0,0)$ is a
  unit timelike vector. Then compute the energy density measured by an observer
  at rest in the coordinates, $w^{\mu} = (1,0,0,0)/\sqrt{1-\beta^{2}}$, where
  $|\beta| < 1$, and compare.
\item \diff{3} \textbf{Designing a fast-clock region.} You want a region where
  clocks run twice as fast as clocks far away, with space kept flat and the
  transition taking place over \SI{10}{\metre}. Using \cref{ex:lapse}, sketch a
  suitable $\alpha(x)$, locate the tensions and pressures it demands, and
  estimate their size in pascals.
\end{problems}

\begin{furtherreading}
\item[\textcite{Schutz2009}] A careful introduction to the stress-energy tensor
  through dust and perfect fluids.
\item[\textcite{Baez2005}] Einstein's equation stated through a ball of
  falling particles, the form used in this chapter, with its consequences
  worked out at an introductory level.
\item[\textcite{Carroll2004}] Einstein's equation, its Newtonian limit and the
  Bianchi identity at graduate level.
\item[\textcite{Hartle2003}] A physical approach to Einstein's equation and to
  what curvature does.
\item[\textcite{Misner1973}] The stress-energy tensor, and why stresses must
  balance in a body at rest.
\item[\textcite{Morris1988a}] The engineer's calculation of this chapter,
  applied to traversable wormholes.
\item[\textcite{Alcubierre1994}] The original warp drive. After \cref{ex:shear}
  you can follow its energy density calculation.
\end{furtherreading}
